Modern Site Reliability Engineering (SRE) teams sit on rich observability data and a deep institutional memory of past incidents recorded in Jira, but the two never meet at decision time. The on-call engineer who is paged at 3 AM still has to manually search Jira for past tickets that look like what they are seeing now --- a slow, brittle process that depends on the engineer remembering the right keywords. The bottleneck in production operations today is no longer ``is something wrong?'' (that is solved by anomaly detectors) but ``what is wrong, has it happened before, and what was the fix?''

This paper describes an end-to-end system that, given a 5-minute window of telemetry (logs, traces, metrics, Kubernetes events), produces three outputs simultaneously: a calibrated triage probability that decides whether the window is worth a ticket at all; a ranked top-5 list of past Jira tickets that look most similar to it; and a flag that fires when no past match exists, marking the failure mode as genuinely new.

Because no public dataset contains rich microservice telemetry paired with the Jira tickets engineers eventually filed for it, we built our own from scratch. We forked Google's Online Boutique microservices demo, added production-realistic OpenTelemetry instrumentation to all ten services under a strict rule (anything a real production service would not emit is forbidden), defined 27 incident scenario families covering hard outages, dependency degradation, pod churn, capacity pressure, near-misses, and network failures, ran 100 independent fault-injection experiments to produce 6{,}720 labeled telemetry windows, and generated a humanized Jira memory corpus of 347 engineer-voice tickets plus 110 distractor tickets designed to look exactly like real tickets but never match anything in the telemetry.

We trained seven independent diagnostic pipelines spanning four model families --- gradient boosting on numeric features, fine-tuned dense bi-encoder retrieval, hybrid sparse$+$dense fusion with knowledge-graph augmentation, log-sequence transformers, Cypher-traversal retrieval over a Neo4j knowledge graph built by an LLM, and a three-stage LLM agent. Every pipeline won on at least one metric and lost on at least one other; no pipeline Pareto-dominated any other. This observation drove the project's central contribution: a four-layer cascade that composes the seven pipelines into a single system. The first layer stacks the six triage scores via class-balanced logistic regression. The second layer fuses four retrievers via Reciprocal Rank Fusion, using a separate multi-retriever overlap-rerank rule to pick the top-1. The third layer detects novel failure modes by OR-combining three independent signals (an LLM agent's verification output, a retrieval-confidence threshold, and a learned per-window classifier). The fourth layer composes the final output.

On 1{,}008 held-out test windows, the cascade achieves Hit@5 $= 0.912$, Hit@1 $= 0.722$, MRR $= 0.794$, strict PR-AUC $= 0.9998$, and a novel-recall of $0.793$ at novel-precision $0.940$ --- a $+388\%$ relative improvement over the previous baseline of $0.163$, and the single largest improvement of the project. The cascade is robust to memory noise (Hit@5 drops only $2\%$ at $50\%$ distractor ratio) and to family-distribution shift (novelty F1 within $11\%$ relative under leave-one-family-out cross-validation). It runs at $16$ microseconds per window at inference --- essentially free.

This report explains every step in plain language: how we injected each kind of fault, what each label means and how it was created, how we used the TAWOS public Jira corpus to ground the humanizer in real engineer voice, how the distractor tickets are designed to confuse the model, what each of the seven models actually is in technical detail, how the four cascade layers feed each other, what experiments validated the design, what the cascade costs to train and deploy, and what evidence we still need to make the result acceptable at a top-tier software engineering venue.
